\section{Conclusion}
\label{sec:conclusion}

We asked whether chain-of-thought reasoning can serve as an interpretability layer for Alzheimer's
biomarker prediction, and we measured the answer rather than asserting it. Across three tasks built
from open cohorts, CoT cost accuracy relative to both black-box tabular baselines
($\Delta$AUC $-0.16$ to $-0.18$; brain-age MAE 11.4\,y vs.\ 5.3\,y) and its own compute-matched
no-reasoning controls, and on the primary amyloid/tau task its chains were demonstrably post-hoc:
the answer was already fixed before the reasoning was written ($\aoc = 0.066$, 82.9\% agreement with
zero chain shown), and corrupting a biomarker value inside the chain barely moved it (3.6\%). A
blinded rubric with a validated shuffled-explanation control nevertheless rated those chains 3.67/4.
That dissociation is the paper's central result, and it makes this a warning rather than a null
result: CoT here produces explanations that would earn clinician trust the model has not justified.

Three qualifications make the picture more than negative. CoT is markedly more faithful on a
continuous target than on classification (\aoc 0.378 vs.\ 0.066), which the multiple-choice-based
faithfulness literature could not have detected. The specific measurements a chain cites do causally
influence its answer (\cfit 1.4--2.1$\times$), even when the surrounding narrative does not. And
both CoT's usefulness and its faithfulness decrease with model capability---it helps and is
load-bearing at 1.5B, and hurts and is decorative at 7B---which inverts the deployment incentive.

The takeaway for practitioners is one sentence: report faithfulness alongside accuracy, or do not
call a reasoning chain an explanation. Four follow-ups are immediate. First, run frontier reasoning
models on the identical harness; this is a drop-in change and the single highest-information
experiment. Second, test the classification-versus-regression faithfulness gap directly by framing
the same underlying quantity as both a binary status and a continuous level. Third, evaluate
constrained-citation prompting---``name the three measurements you used''---which \cfit suggests may
retain most of the interpretability value without the accuracy cost of free-form narration. Fourth,
build the hybrid: a tabular foundation model for the prediction and an LLM for the narration,
conditioned on the model's attributions and scored by this battery, so the explanation is causally
anchored by construction.
